\IfFileExists{stacks-project.cls}{%
\documentclass{stacks-project}
}{%
\documentclass{amsart}
}

% For dealing with references we use the comment environment
\usepackage{verbatim}
\newenvironment{reference}{\comment}{\endcomment}
%\newenvironment{reference}{}{}
\newenvironment{slogan}{\comment}{\endcomment}
\newenvironment{history}{\comment}{\endcomment}

% For commutative diagrams we use Xy-pic
\usepackage[all]{xy}

% We use 2cell for 2-commutative diagrams.
\xyoption{2cell}
\UseAllTwocells

% We use multicol for the list of chapters between chapters
\usepackage{multicol}

% This is generally recommended for better output
\usepackage{lmodern}
\usepackage[T1]{fontenc}

% For cross-file-references
\usepackage{xr-hyper}

% Package for hypertext links:
\usepackage{hyperref}

% For any local file, say "hello.tex" you want to link to please
% use \externaldocument[hello-]{hello}
\externaldocument[introduction-]{introduction}
\externaldocument[conventions-]{conventions}
\externaldocument[sets-]{sets}
\externaldocument[categories-]{categories}
\externaldocument[topology-]{topology}
\externaldocument[sheaves-]{sheaves}
\externaldocument[sites-]{sites}
\externaldocument[stacks-]{stacks}
\externaldocument[fields-]{fields}
\externaldocument[algebra-]{algebra}
\externaldocument[brauer-]{brauer}
\externaldocument[homology-]{homology}
\externaldocument[derived-]{derived}
\externaldocument[simplicial-]{simplicial}
\externaldocument[more-algebra-]{more-algebra}
\externaldocument[smoothing-]{smoothing}
\externaldocument[modules-]{modules}
\externaldocument[sites-modules-]{sites-modules}
\externaldocument[injectives-]{injectives}
\externaldocument[cohomology-]{cohomology}
\externaldocument[sites-cohomology-]{sites-cohomology}
\externaldocument[dga-]{dga}
\externaldocument[dpa-]{dpa}
\externaldocument[sdga-]{sdga}
\externaldocument[hypercovering-]{hypercovering}
\externaldocument[schemes-]{schemes}
\externaldocument[constructions-]{constructions}
\externaldocument[properties-]{properties}
\externaldocument[morphisms-]{morphisms}
\externaldocument[coherent-]{coherent}
\externaldocument[divisors-]{divisors}
\externaldocument[limits-]{limits}
\externaldocument[varieties-]{varieties}
\externaldocument[topologies-]{topologies}
\externaldocument[descent-]{descent}
\externaldocument[perfect-]{perfect}
\externaldocument[more-morphisms-]{more-morphisms}
\externaldocument[flat-]{flat}
\externaldocument[groupoids-]{groupoids}
\externaldocument[more-groupoids-]{more-groupoids}
\externaldocument[etale-]{etale}
\externaldocument[chow-]{chow}
\externaldocument[intersection-]{intersection}
\externaldocument[pic-]{pic}
\externaldocument[weil-]{weil}
\externaldocument[adequate-]{adequate}
\externaldocument[dualizing-]{dualizing}
\externaldocument[duality-]{duality}
\externaldocument[discriminant-]{discriminant}
\externaldocument[derham-]{derham}
\externaldocument[local-cohomology-]{local-cohomology}
\externaldocument[algebraization-]{algebraization}
\externaldocument[curves-]{curves}
\externaldocument[resolve-]{resolve}
\externaldocument[models-]{models}
\externaldocument[functors-]{functors}
\externaldocument[equiv-]{equiv}
\externaldocument[pione-]{pione}
\externaldocument[etale-cohomology-]{etale-cohomology}
\externaldocument[proetale-]{proetale}
\externaldocument[relative-cycles-]{relative-cycles}
\externaldocument[more-etale-]{more-etale}
\externaldocument[trace-]{trace}
\externaldocument[crystalline-]{crystalline}
\externaldocument[spaces-]{spaces}
\externaldocument[spaces-properties-]{spaces-properties}
\externaldocument[spaces-morphisms-]{spaces-morphisms}
\externaldocument[decent-spaces-]{decent-spaces}
\externaldocument[spaces-cohomology-]{spaces-cohomology}
\externaldocument[spaces-limits-]{spaces-limits}
\externaldocument[spaces-divisors-]{spaces-divisors}
\externaldocument[spaces-over-fields-]{spaces-over-fields}
\externaldocument[spaces-topologies-]{spaces-topologies}
\externaldocument[spaces-descent-]{spaces-descent}
\externaldocument[spaces-perfect-]{spaces-perfect}
\externaldocument[spaces-more-morphisms-]{spaces-more-morphisms}
\externaldocument[spaces-flat-]{spaces-flat}
\externaldocument[spaces-groupoids-]{spaces-groupoids}
\externaldocument[spaces-more-groupoids-]{spaces-more-groupoids}
\externaldocument[bootstrap-]{bootstrap}
\externaldocument[spaces-pushouts-]{spaces-pushouts}
\externaldocument[spaces-chow-]{spaces-chow}
\externaldocument[groupoids-quotients-]{groupoids-quotients}
\externaldocument[spaces-more-cohomology-]{spaces-more-cohomology}
\externaldocument[spaces-simplicial-]{spaces-simplicial}
\externaldocument[spaces-duality-]{spaces-duality}
\externaldocument[formal-spaces-]{formal-spaces}
\externaldocument[restricted-]{restricted}
\externaldocument[spaces-resolve-]{spaces-resolve}
\externaldocument[formal-defos-]{formal-defos}
\externaldocument[defos-]{defos}
\externaldocument[cotangent-]{cotangent}
\externaldocument[examples-defos-]{examples-defos}
\externaldocument[algebraic-]{algebraic}
\externaldocument[examples-stacks-]{examples-stacks}
\externaldocument[stacks-sheaves-]{stacks-sheaves}
\externaldocument[criteria-]{criteria}
\externaldocument[artin-]{artin}
\externaldocument[quot-]{quot}
\externaldocument[stacks-properties-]{stacks-properties}
\externaldocument[stacks-morphisms-]{stacks-morphisms}
\externaldocument[stacks-limits-]{stacks-limits}
\externaldocument[stacks-cohomology-]{stacks-cohomology}
\externaldocument[stacks-perfect-]{stacks-perfect}
\externaldocument[stacks-introduction-]{stacks-introduction}
\externaldocument[stacks-more-morphisms-]{stacks-more-morphisms}
\externaldocument[stacks-geometry-]{stacks-geometry}
\externaldocument[moduli-]{moduli}
\externaldocument[moduli-curves-]{moduli-curves}
\externaldocument[examples-]{examples}
\externaldocument[exercises-]{exercises}
\externaldocument[guide-]{guide}
\externaldocument[desirables-]{desirables}
\externaldocument[coding-]{coding}
\externaldocument[obsolete-]{obsolete}
\externaldocument[fdl-]{fdl}
\externaldocument[index-]{index}

% Theorem environments.
%
\theoremstyle{plain}
\newtheorem{theorem}[subsection]{Theorem}
\newtheorem{proposition}[subsection]{Proposition}
\newtheorem{lemma}[subsection]{Lemma}

\theoremstyle{definition}
\newtheorem{definition}[subsection]{Definition}
\newtheorem{example}[subsection]{Example}
\newtheorem{exercise}[subsection]{Exercise}
\newtheorem{situation}[subsection]{Situation}

\theoremstyle{remark}
\newtheorem{remark}[subsection]{Remark}
\newtheorem{remarks}[subsection]{Remarks}

\numberwithin{equation}{subsection}

% Macros
%
\def\lim{\mathop{\mathrm{lim}}\nolimits}
\def\colim{\mathop{\mathrm{colim}}\nolimits}
\def\Spec{\mathop{\mathrm{Spec}}}
\def\Hom{\mathop{\mathrm{Hom}}\nolimits}
\def\Ext{\mathop{\mathrm{Ext}}\nolimits}
\def\SheafHom{\mathop{\mathcal{H}\!\mathit{om}}\nolimits}
\def\SheafExt{\mathop{\mathcal{E}\!\mathit{xt}}\nolimits}
\def\Sch{\mathit{Sch}}
\def\Mor{\mathop{\mathrm{Mor}}\nolimits}
\def\Ob{\mathop{\mathrm{Ob}}\nolimits}
\def\Sh{\mathop{\mathit{Sh}}\nolimits}
\def\NL{\mathop{N\!L}\nolimits}
\def\CH{\mathop{\mathrm{CH}}\nolimits}
\def\proetale{{pro\text{-}\acute{e}tale}}
\def\etale{{\acute{e}tale}}
\def\QCoh{\mathit{QCoh}}
\def\Ker{\mathop{\mathrm{Ker}}}
\def\Im{\mathop{\mathrm{Im}}}
\def\Coker{\mathop{\mathrm{Coker}}}
\def\Coim{\mathop{\mathrm{Coim}}}

% Boxtimes
%
\DeclareMathSymbol{\boxtimes}{\mathbin}{AMSa}{"02}

%
% Macros for moduli stacks/spaces
%
\def\QCohstack{\mathcal{QC}\!\mathit{oh}}
\def\Cohstack{\mathcal{C}\!\mathit{oh}}
\def\Spacesstack{\mathcal{S}\!\mathit{paces}}
\def\Quotfunctor{\mathrm{Quot}}
\def\Hilbfunctor{\mathrm{Hilb}}
\def\Curvesstack{\mathcal{C}\!\mathit{urves}}
\def\Polarizedstack{\mathcal{P}\!\mathit{olarized}}
\def\Complexesstack{\mathcal{C}\!\mathit{omplexes}}
% \Pic is the operator that assigns to X its picard group, usage \Pic(X)
% \Picardstack_{X/B} denotes the Picard stack of X over B
% \Picardfunctor_{X/B} denotes the Picard functor of X over B
\def\Pic{\mathop{\mathrm{Pic}}\nolimits}
\def\Picardstack{\mathcal{P}\!\mathit{ic}}
\def\Picardfunctor{\mathrm{Pic}}
\def\Deformationcategory{\mathcal{D}\!\mathit{ef}}
\usepackage{cleveref}
\crefname{originallemma}{original lemma}{original lemmas}
\Crefname{originallemma}{Original lemma}{Original lemmas}
\makeatletter
\newlabel{lemma-groupoid-on-field-geometrically-irreducible}{{04LT}{source}{}{sourcepointer-1}{}}
\newlabel{lemma-groupoid-on-field-geometrically-irreducible@cref}{{[originallemma][0][]04LT}{[1][0][]source}{}{}{}}
\makeatother
\crefname{originallemma}{original lemma}{original lemmas}
\Crefname{originallemma}{Original lemma}{Original lemmas}
\makeatletter
\newlabel{varieties-lemma-integral-closure-ground-field}{{04MI}{source}{}{sourcepointer-2}{}}
\newlabel{varieties-lemma-integral-closure-ground-field@cref}{{[originallemma][0][]04MI}{[1][0][]source}{}{}{}}
\makeatother
\crefname{originalproposition}{original proposition}{original propositions}
\Crefname{originalproposition}{Original proposition}{Original propositions}
\makeatletter
\newlabel{varieties-proposition-unique-base-field}{{04MK}{source}{}{sourcepointer-3}{}}
\newlabel{varieties-proposition-unique-base-field@cref}{{[originalproposition][0][]04MK}{[1][0][]source}{}{}{}}
\makeatother
\crefname{originallemma}{original lemma}{original lemmas}
\Crefname{originallemma}{Original lemma}{Original lemmas}
\makeatletter
\newlabel{lemma-groupoid-on-field-locally-finite-type-dimension}{{04LU}{source}{}{sourcepointer-4}{}}
\newlabel{lemma-groupoid-on-field-locally-finite-type-dimension@cref}{{[originallemma][0][]04LU}{[1][0][]source}{}{}{}}
\makeatother
\crefname{originallemma}{original lemma}{original lemmas}
\Crefname{originallemma}{Original lemma}{Original lemmas}
\makeatletter
\newlabel{groupoids-lemma-group-scheme-finite-type-field}{{045X}{source}{}{sourcepointer-5}{}}
\newlabel{groupoids-lemma-group-scheme-finite-type-field@cref}{{[originallemma][0][]045X}{[1][0][]source}{}{}{}}
\makeatother
\title{Dimension of a groupoid scheme over a field with locally finite type structure maps equals the dimension of its stabilizer group scheme}
\author{The Stacks Project Authors}
\date{stacks-project:more-groupoids.tex:local-snapshot-20261012; GNU FDL 1.2 or later}
\begin{document}
\maketitle
\setcounter{subsection}{1}
\noindent\textit{Editorial note.} All delivered mathematical text is the author's own, copied from the pinned Stacks Project chapter file more-groupoids.tex; no statement or proof step was rewritten, reordered or removed. The counted span is the author environment and its own proof environment. References inside the span are handed to the assembler as navigation stubs that carry the canonical tag of the cited tag only; whether each of them is load-bearing is recorded on the stub itself.\par\medskip
\noindent\textit{Allowed standard prerequisites (editorial).}\begin{itemize}
\item Uniqueness and geometric irreducibility of the irreducible component through the identity for a groupoid over a field (more-groupoids chapter, Lemma lemma-groupoid-on-field-geometrically-irreducible, canonical tag 04LT; collected on this lane as L22510). cited at proof lines 1562-1564. Provides the component Z through e, unique, thought of as an integral closed subscheme of R.
\item The integral closure of the ground field in a function field of a variety is a field (varieties chapter, Lemma varieties-lemma-integral-closure-ground-field, canonical tag 04MI). cited at proof lines 1565-1568. Guarantees k'\_s and k'\_t are fields, the form in which the unique base field proposition applies.
\item Uniqueness of the base field inside the function field of an integral variety (varieties chapter, Proposition varieties-proposition-unique-base-field, canonical tag 04MK). cited at proof lines 1569-1571. Gives k'\_s = k'\_t as subrings of Gamma(Z, O\_Z), the pivot for s(k) = t(k).
\item Equidimensionality of a groupoid over a field with locally finite type structure maps (more-groupoids chapter, Lemma lemma-groupoid-on-field-locally-finite-type-dimension, canonical tag 04LU; also in this lane's queue). cited at proof lines 1584-1586. Gives dim(R) = dim(Z), one side of the dimension squeeze.
\item Equidimensionality of group schemes of finite type over a field (groupoids chapter, Lemma groupoids-lemma-group-scheme-finite-type-field, canonical tag 045X). cited at proof lines 1584-1586. Gives the equidimensionality of G, the other side of the squeeze.
\end{itemize}\medskip
\ifdefined\hypertarget\hypertarget{sourcepointer-1}{}\fi
\noindent\textit{Source reference.} Original lemma 04LT: sources/stacks-chow/tags.txt, source lines 6019--6019. Cited at proof lines 1562-1564: uniqueness of the irreducible component Z through e.\par\medskip
\ifdefined\hypertarget\hypertarget{sourcepointer-2}{}\fi
\noindent\textit{Source reference.} Original lemma 04MI: sources/stacks-chow/tags.txt, source lines 6044--6044. Cited at proof lines 1565-1568: k'\_s and k'\_t are fields.\par\medskip
\ifdefined\hypertarget\hypertarget{sourcepointer-3}{}\fi
\noindent\textit{Source reference.} Original proposition 04MK: sources/stacks-chow/tags.txt, source lines 6046--6046. Cited at proof lines 1569-1571: k'\_s = k'\_t as subrings of Gamma(Z, O\_Z).\par\medskip
\ifdefined\hypertarget\hypertarget{sourcepointer-4}{}\fi
\noindent\textit{Source reference.} Original lemma 04LU: sources/stacks-chow/tags.txt, source lines 6020--6020. Cited at proof lines 1584-1586: R is equidimensional, giving dim(R) = dim(Z).\par\medskip
\ifdefined\hypertarget\hypertarget{sourcepointer-5}{}\fi
\noindent\textit{Source reference.} Original lemma 045X: sources/stacks-chow/tags.txt, source lines 5452--5452. Cited at proof lines 1584-1586: G is equidimensional, the other side of the squeeze.\par\medskip
\begin{lemma}
\label{lemma-groupoid-on-field-dimension-equal-stabilizer}
Let $S$ be a scheme. Let $(U, R, s, t, c)$ be a groupoid scheme
over $S$. Assume $U = \Spec(k)$ with $k$ a field.
Assume $s, t$ are locally of finite type.
Then $\dim(R) = \dim(G)$ where $G$ is the stabilizer group scheme of $R$.
\end{lemma}

\begin{proof}
Let $Z \subset R$ be the irreducible component passing through $e$ (see
Lemma \ref{lemma-groupoid-on-field-geometrically-irreducible})
thought of as an integral closed subscheme of $R$.
Let $k'_s$, resp.\ $k'_t$ be the integral closure of
$s(k)$, resp.\ $t(k)$ in $\Gamma(Z, \mathcal{O}_Z)$.
Recall that $k'_s$ and $k'_t$ are fields, see
Varieties, Lemma \ref{varieties-lemma-integral-closure-ground-field}.
By
Varieties, Proposition \ref{varieties-proposition-unique-base-field}
we have $k'_s = k'_t$ as subrings of $\Gamma(Z, \mathcal{O}_Z)$.
As $e$ factors through $Z$ we obtain a commutative diagram
$$
\xymatrix{
k \ar[rd]_t \ar[rrd]^1 \\
& \Gamma(Z, \mathcal{O}_Z) \ar[r]^e & k \\
k \ar[ru]^s \ar[rru]_1
}
$$
This on the one hand shows that $k'_s = s(k)$, $k'_t = t(k)$, so
$s(k) = t(k)$, which combined with the diagram above implies
that $s = t$! In other words, we conclude that $Z$ is a closed
subscheme of $G = R \times_{(t, s), U \times_S U, \Delta} U$.
The lemma follows as both $G$ and $R$ are equidimensional, see
Lemma \ref{lemma-groupoid-on-field-locally-finite-type-dimension} and
Groupoids, Lemma \ref{groupoids-lemma-group-scheme-finite-type-field}.
\end{proof}

\end{document}
